
\documentclass[oneside,12pt,numbers,spanish]{ezthesis}
%% # Opciones disponibles para el documento #
%%
%% Las opciones con un (*) son las opciones predeterminadas.
%%
%% Modo de compilar:
%%   draft            - borrador con marcas de fecha y sin im'agenes
%%   draftmarks       - borrador con marcas de fecha y con im'agenes
%%   final (*)        - version final de la tesis
%%
%% Tama'no de papel:
%%   letterpaper (*)  - tama'no carta (Am'erica)
 %%   a4paper          - tama'no A4    (Europa)
%%
%% Formato de impresi'on:
%%   oneside          - hojas impresas por un solo lado
%%   twoside (*)      - hijas impresas por ambos lados
%%
%% Tama'no de letra:
%%   10pt, 11pt, o 12pt (*)
%%
%% Espaciado entre renglones:
%%   singlespace      - espacio sencillo
%%   onehalfspace (*) - espacio de 1.5
%%   doublespace      - a doble espacio
%%
%% Formato de las referencias bibliogr'aficas:
%%   numbers          - numeradas, p.e. [1]
%%   authoryear (*)   - por autor y a'no, p.e. (Newton, 1997)
%%
%% Opciones adicionales:
%%   spanish         - tesis escrita en espa'nol
%%
%% Desactivar opciones especiales:
%%   nobibtoc   - no incluir la bibiolgraf'ia en el 'Indice general
%%   nofancyhdr - no incluir "fancyhdr" para producir los encabezados
%%   nocolors   - no incluir "xcolor" para producir ligas con colores
%%   nographicx - no incluir "graphicx" para insertar gr'aficos
%%   nonatbib   - no incluir "natbib" para administrar la bibliograf'ia

%% Paquetes adicionales requeridos se pueden agregar tambi'en aqu'i.
%% Por ejemplo:
\usepackage[spanish]{babel}
\usepackage{cite}
\usepackage{natbib} 
\usepackage{subfig}
\usepackage[utf8x]{inputenc}
\usepackage{multicol}
\usepackage{multirow, array}
\usepackage{graphicx} 
\usepackage{float}  
\usepackage{amsmath}
\usepackage{amssymb}
\usepackage{Pgfgantt }
\pretolerance = 2000
\tolerance = 3000
\righthyphenmin = 2000
\lefthyphenmin = 2000


%% # Datos del documento #
%%\documentclass[12pt]{•}%% Nota que los acentos se deben escribir: \'a, \'e, \'i, etc.
%% La letra n con tilde es: \~n.

\author{Estefany C. Salas S.}
\title{Clasificación de tubérculos de papa criolla para diferentes densidades de siembra según el peso fresco por calibre empleando Bayes Naive.}
\supervisor{ Rossana Timaure}
\institution{Universidad Nacional Experimental del T\'achira}
\faculty{Vicerectorado Acad\'emico}
\department{Departamento de Ingenier\'ia en Inform\'atica}

%% # M'argenes del documento #
%% 
%% Quitar el comentario en la siguiente linea para austar los m'argenes del
%% documento. Leer la documentaci'on de "geometry" para m'as informaci'on.

\geometry{left=2cm,right=2cm,top=3cm,bottom=3cm}

%% El siguiente comando agrega ligas activas en el documento para las
%% referencias cruzadas y citas bibliogr'aficas. Tiene que ser *la 'ultima*
%% instrucci'on antes de \begin{document}.
\hyperlinking
\begin{document}

%% En esta secci'on se describe la estructura del documento de la tesis.
%% Consulta los reglamentos de tu universidad para determinar el orden
%% y la cantidad de secciones que debes de incluir.

%% # Portada de la tesis #
%% Mirar el archivo "titlepage.tex" para los detalles.
%% ## Construye tu propia portada ##
%% 
%% Una portada se conforma por una secuencia de "Blocks" que incluyen
%% piezas individuales de informaci'on. Un "Block" puede incluir, por
%% ejemplo, el t'itulo del documento, una im'agen (logotipo de la universidad),
%% el nombre del autor, nombre del supervisor, u cualquier otra pieza de
%% informaci'on.
%%
%% Cada "Block" aparece centrado horizontalmente en la p'agina y,
%% verticalmente, todos los "Blocks" se distruyen de manera uniforme 
%% a lo largo de p'agina.
%%
%% Nota tambi'en que, dentro de un mismo "Block" se pueden cortar
%% lineas usando el comando \\
%%
%% El tama'no del texto dentro de un "Block" se puede modificar usando uno de
%% los comandos:
%%   \small      \LARGE
%%   \large      \huge
%%   \Large      \Huge
%%
%% Y el tipo de letra se puede modificar usando:
%%   \bfseries - negritas
%%   \itshape  - it'alicas
%%   \scshape  - small caps
%%   \slshape  - slanted
%%   \sffamily - sans serif
%%
%% Para producir plantillas generales, la informaci'on que ha sido inclu'ida
%% en el archivo principal "tesis.tex" se puede accesar aqu'i usando:
%%   \insertauthor
%%   \inserttitle
%%   \insertsupervisor
%%   \insertinstitution
%%   \insertdegree
%%   \insertfaculty
%%   \insertdepartment
%%   \insertsubmitdate

\begin{titlepage}
  \TitleBlock{\includegraphics[scale=0.2]{unet.jpg} }  
  \TitleBlock{\small\insertinstitution}
  \TitleBlock[\small\bigskip]{\insertfaculty}
  \TitleBlock[\small\bigskip]{Decanato de Docencia}
  \TitleBlock[\small\bigskip]{\insertdepartment}
  \TitleBlock[\small\bigskip]{Trabajo de aplicaci\'on profesional}
  \TitleBlock[\small\bigskip]{Proyecto especial de grado}
    
 \TitleBlock{\scshape\inserttitle}
 
 \begin{flushright}
  \TitleBlock{ \small Autor(es):  \insertauthor} 
  \TitleBlock{ \small C.I.: 24820631}
	\TitleBlock{ \small estefany.salas@unet.edu.ve}
  \TitleBlock{ \small }
  \TitleBlock{\small Tutor(es):  \insertsupervisor}
  \TitleBlock{ \small rttg@unet.edu.ve}
 \end{flushright}
  \TitleBlock{\small San Crist\'obal, Julio de 2018.}
  
\end{titlepage}

%% Nota 1:
%% Se puede agregar un escudo o logotipo en un "Block" como:
%%   \TitleBlock{\includegraphics[height=4cm]{escudo_uni}}
%% y teniendo un archivo "escudo_uni.pdf", "escudo_uni.png" o "escudo_uni.jpg"
%% en alg'un lugar donde LaTeX lo pueda encontrar.

%% Nota 2:
%% Normalmente, el espacio entre "Blocks" se extiende de modo que el
%% contenido se reparte uniformemente sobre toda la p'agina. Este
%% comportamiento se puede modificar para mantener fijo, por ejemplo, el
%% espacio entre un par de "Blocks". Escribiendo:
%%   \TitleBlock{Bloque 1}
%%   \TitleBlock[\bigskip]{Bloque2}
%% se deja un espacio "grande" y de tama~no fijo entre el bloque 1 y 2.
%% Adem'as de \bigskip est'an tambi'en \smallskip y \medskip. Si necesitas
%% aun m'as control puedes usar tambi'en, por ejemplo, \vspace*{2cm}.




% \tableofcontents
% \listoffigures

\righthyphenmin = 2000
\lefthyphenmin = 2000

\chapter*{Introducci\'on}
\pagenumbering{arabic} % para empezar la numeración con números

En Colombia, el nombre papa criolla corresponde a un morfotipo que desarrolla tubérculos con piel y pulpa de color amarillo, esta variedad ha sido clasificada como \textit{Solanum phureja}, según Fedepapa \textcolor{red}{ QUE ES FEDEPAPA} esta variedad de papa es un alimento de alto valor nutritivo y de excelentes calidades culinarias, además de ser una de las fuentes de proteína más económica y ser considerada un producto exótico por consumidores de Europa y Estados Unidos, lo que le permite posicionarse como un producto para exportación que normalmente se comercializa en presentaciones que van desde paquetes de papas congeladas a latas de papa.\\

Al ser un producto de alta demanda requiere altos controles de calidad  y un  manejo especializado de las variables que influyen en su crecimiento, por supuesto las industrias que trabajan con este tubérculo buscan opciones que les permitan mejorar la calidad de la papa incurriendo en los menores gastos, aquí entran variables como las que van a ser investigadas por el presente, que son la densidad de siembra (distancia entre plantas y distancia entre surcos) y los pesos frescos de cada calibre, en investigaciones anteriores se ha determinado que existe una relación entre estas variables, dichas investigaciones señalan que la densidad de siembra tiene un efecto sobre los tamaños y cantidad de tubérculos que la misma va a producir, esto indicaría que esta variable podría ser usada por industrias para obtener la mayor cantidad de tubérculos que se adecuen a los tamaños que ellos desean. \\

Esta investigación tiene como objetivo crear un clasificador empleando \textit{Bayes Naive} que teniendo como variables de entrada (características) los pesos frescos de cada calibre, logre identificar con un nivel alto de aciertos la densidad de siembra que fue usada en esa cosecha, y así poder establecer la relacion entre estás variables de manera más precisa. Se usará como método de clasificación \textit{Bayes Naive} porque ha demostrado dar un gran porcentaje de aciertos al momento de clasificar y no posee limitaciones como los métodos estadísticos clásicos o la regresión espacial. Los datos observados para realizar el modelo de clasificación fueron recolectados previamente por un estudiante de la Universidad Nacional de Colombia, a partir de estos se plantea generar un modelo empleando regresión espacial con el cual se pueda generar datos artificiales, para luego realizar el clasificador con ambos conjuntos de datos y evaluar cual se ajusta mejor, evaluando la curva característica operativa del receptor para cada caso.\\ 

La herramienta informática a emplear para desarrollar el algoritmo y realizar la clasificación es Python. Python es un lenguaje de programación reconocido por poseer una sintaxis simple y ser fácil de aprender, además es eficiente, este lenguaje permite a más personas aprender a programar de manera sencilla y así concentrarse en los problemas que buscan resolver. Actualmente se puede observar una gran tendencia al uso de Python en grandes centros de investigación como el CERN (Organización Europea para la Investigación Nuclear) y por parte de científicos en ramas como la Bioinformática, Neurofisiología, Física, Matemáticas, etc. Esto es debido a la disponibilidad de librerías de visualización, procesamiento de señales, estadísticas, álgebra, etc.; de fácil utilización y que cuentan con muy buena documentación. Los paquetes de Python denominados NumPy (Python Númerico) y SciPy (Python Científico) son pilares para la realización de trabajos científicos capaces de emular las funciones de otros lenguajes netamente científicos con Matlab.\\


 \textcolor{red}{ TODO DEBE SER REDACTADO EN PASADO PERFECTO}


 \textcolor{red}{ Se debe agregar esto modificar la redaccion a gusto y completar el contenido}

Este documento consta de cinco capítulos en los cuales se describe el desarrollo  de la investigación  y estan estructurados de la siguiente forma:\\
Capitulo 1. Preliminares: Donde se describe el planteamiento del problema, el objetivo general planteado y los objetivos específicos que llevaron  a cabo para cumplir con el objetivo general de  la investigación.\\
Capítulo 2. Marco teórico: Contiene los antecedentes y la perspectiva teórica, para el entendimiendo
de los conceptos relacionados con esta investigación.\\
Capítulo 3. Metodología: Describe la metodología utilizada durante la investigación, el enfoque,
tipo, nivel y dise\~no de la misma y la metodología con la que se desarrolló el algoritmo.\\
Capítulo 4. \\
Capítulo 5.\\


%En Colombia, el nombre papa criolla corresponde a un morfotipo que desarrolla tubérculos con piel color amarilla y pulpa ("`yema de huevo"')(Rodriguez, 2009), esta variedad ha sido clasificada como \textit{Solanum phureja} (Hawkes, 1990), según Fedepapa en 1988 esta variedad de papa es un alimento de alto valor nutritivo y de excelentes calidades culinarias, además de ser una de las fuentes de proteína más económica y ser considerada un producto exótico por consumidores de Europa y Estados Unidos, lo que le permite posicionarse como un producto para exportación que normalmente se comercializa en presentaciones que van desde paquetes de papas congeladas a latas de papa. Según Rivera et al (2011) los tubérculos de papa criolla como cualquier otro producto deben satisfacer requerimientos de calidad externos e internos que son impuestos por el uso que se le va a dar, por ejemplo, algunos de los requerimientos que caracterizan la calidad culinaria de la papa son la textura, el aroma, el sabor y el color, pero hay muchas más caracteristicas observables como el tamaño el cual es un factor clave en el uso industrial de la papa, Piñeros (2009) indica que los tamaños pequeños son apropiados para las presentaciones precocida congelada o encurtida y los tamaños grandes con formas redondas o comprimidas son apropiadas en procesos de obtención de francesas y hojuelas.\\ 

%Cuando se toma en cuenta la importancia del tamaño de la papa para la industria se empiezan a observar investigaciones que buscan afectar la variable tamaño con el menor costo posible, es decir, evitando químicos y métodos artificiales para el control del crecimiento, algunos investigadores como Arias et al (1996) pusieron a evaluacion el peso y número de tubérculos de cada clase de papa para diferentes densidades de siembra, por otra parte Bernal (2017) empleó modelos de regresión binomial negativa cero-inflada que evidenciaron el efecto significativo de la densidad de siembra sobre el conteo de tubérculos en los calibres superiores a 4 cm mientras que el modelo binomial negativo lo evidenció en el caso de los calibres inferiores a 4 cm. Las investigaciones toman en cuenta la variable densidad de siembra para evaluarla como factor influyente en el tamaño de los tubérculos porque es una variable controlada, que puede ser modificada sin alterar abruptamente los costos de la siembra, sin embargo, esta variable que representa la distancia entre plantas y surcos posee caracteristicas espaciales.\\

%Al realizar clasificaciones de la densidad de siembra en relación con otras variables como el peso no se pueden implementar métodos estadisticos como la regresión o el análisis de varianza porque se incumplen los supuestos para la aplicación de estos al estar frente a una variable espacial, por lo tanto, el método estadístico correcto sería la regresión espacial, aunque al aplicar este método existe otro problema y es que la variable de densidad de siembra a pesar de tener varios valores es fija, por lo tanto se debe realizar el análisis para cada valor de densidad, pero realizar esto no asegura una corelacción general de las densidades con el tamaño de los tubérculos aunque cada análisis individual de un resultado similar, esto lleva a buscar soluciones alternativas como la que va a ser evaluada en esta investigación que es la clasificación empleando el algoritmo de Bayes Naïve que según la descripción de Misigo y Miriti (2016) es un clasificador probabilístico basado en el teorema de Bayes considerando la suposición de independencia de Naive, este clasificador asume que el valor de una variable en una clase es independiente del valor de otra variable, con Bayes Naïve se pueden llevar a cabo métodos de clasificación más sofisticados que permiten crear modelos con capacidades de predición.


\include{Capitulo1}
\include{Capitulo2}
\include{Capitulo31}
%% Los cap'itulos inician con \chapter{T'itulo}, estos aparecen numerados y
%% se incluyen en el 'indice general.
%%
%% Recuerda que aqu'i ya puedes escribir acentos como: 'a, 'e, 'i, etc.
%% La letra n con tilde es: 'n.

\chapter*{Referencias Bibliográficas}

\noindent
Althoff, C. (2016). "`The Self-taught Programmer"'. Cory Althoff.\\

\noindent
Arbia, G. (2014). "`A Primer for Spatial Econometrics With Applications in R"'. Palgrave Macmillan.\\

\noindent
Arias, V., Bustos, P., Ñustéz, C. (1996). "`Evaluación del Rendimiento en papa criolla (Solanum phureja) variedad "`Yema de Huevo"', bajo diferentes Densidades de Siembra en la Sabana de Bogotá."' Agronomía Colombiana, Vol. XIII, No 2, pp. 152-161.\\

\noindent
Bernal, N., Darghan, AE., Rodríguez, LE. (2017). "`Modelado del Calibre de tubérculos de papa Solanum phureja bajo diferentes densidades de siembra mediante regresión binomial negativa cero-inflada"'.\\

\noindent
Bonaccorso, G. (2017).  "`Machine Learning Algorithms."'. Packt Publishing Ltd.

\noindent
Buitrago, G., López, A., Coronado, A., Osorno, F. (2004). "`Determinación de las características físicas y propiedades mecánicas de papa cultivada en Colombia"'. Revista Brasileira de Engenharia Agrícola e Ambiental, Vol.8, No.1, pp.102-110.\\

\noindent
Challenger, I., Díaz, Y., Becerra, R. (2014). "`El lenguaje de programación Python"'. Ciencias Holguín, Vol. XX, No. 2, pp. 1-13.\\

\noindent
Gönen, M. (2007). "`Analyzing Receiver Operating Characteristic Curves with SAS."' SAS Institute Inc.\\

\noindent
IBM. (2012). "`Manual CRISP-DM de IBM SPSS Modeler"'. IBM Corporation.\\

\noindent
Koduvely, H. (2015). "`Learning Bayesian Models with R"'. Reino Unido: Packt Publishing Ltd.\\

\noindent
Layton, R. (2015). "`Learning Data Mining with Python"'. Packt Publishing.\\

\noindent
Misigo, R., Miriti, E. (2016). "`Classification of Selected Apple Fruit Varieties using Naive Bayes"'. Indian Journal of Computer Science and Engineering (IJCSE), Vol. 7, No. 1.\\

\noindent
Mohamad, N., Jusoh, N., Htike, Z., Lei Win, S. (2014). "`Bacteria Identification from Microscopic Morphology Using Naïve Bayes."' International Journal of Computer Science, Engineering and Information Technology (IJCSEIT), Vol. 4, No.2.\\

\noindent 
Piñeros, C. (2009). "`Recopilación de la Investigación del Sistema Productivo Papa Criolla"'. Secretaria de Agricultura y Desarrollo Economico (Gobernación de Cundinamarca), Federación Colombiana de Productos de Papa.\\

\noindent
Rivera, J., Herrera, A., Rodríguez, L. (2011). "`Evaluación de la aptitud de procesamiento en seis genotipos de papa criolla (Solanum tuberosum Grupo Phureja)"'. Agronomía Colombiana, Vol. 29(1), pp. 73-81.\\

\noindent
Sampieri, R., Fernández, C., Baptista, M. (2014). "`Metodología de la Investigación"'. McGRAW-HILL / INTERAMERICANA EDITORES, S.A. DE C.V.\\

\noindent
Tsangaratos, P., Ilia, I. (2016). "`Comparison of a logistic regression and Naïve Bayes classifier in landslide susceptibility assessments: The influence of models complexity and training dataset size"'. Catena 145, pp. 164–179.\\

\noindent
VanderPlas, J. (2017). "`Python Data Science Handbook"'. O’Reilly.\\


\noindent
% Aguilar, Ma.Guadalupe, José Carrillo, Antonio Rivera, and Victor González. 2006. “Growth Analysis and Sink-Source Relationships in Two Potato (Solanum tuberosum L.) Varieties.” Revista Fitotecnia Mexicana 29(2):145–56.\\

\noindent
%Available CRAN Packages By Name. Disponible en:\url{ https://cran.r-project.org/web/packages/available_packages_by_name.html}. Consultada Enero 2018.\\

\noindent
% Barraza, F.V., Fischer, G.,Cardona, C.E. (2004). Studying the process of tomato crop (Lycopersicon esculentum Mill.) growth in the
%Middle Sinu Valley, Colombia. Agronomía Colombiana 22 (1): 81-90 \\

\noindent
% Blackman, V.H. (1920). The Significance of the Efficiency Index of Plant Growth. New Phytologist, Vol. 19, No. 3/4, pp. 97-100\\

\noindent
% Borrego, Fernando et al. 2000. “Nota Técnica. Anáisis de Crecimiento En Siete Variedades de Papa (Solanum Tuberosum L.).” Agronomía Mesoaméricana 11(1):145-49.\\

\noindent
%Chambers, John M (1998). Programming with Data: A Guide to the S Language. Springer.\\

\noindent
% Condori, Bruno et al. 2008. “Agrophysiological Characterisation and Parametrisation of Andean Tubers: Potato (Solanum sp.), Oca (Oxalis tuberosa), Isaño (Tropaeolum tuberosum) and Papalisa (Ullucus tuberosus).” European Journal of Agronomy 28(4):526-40.\\

\noindent
% De Mendiburu F. (2017).  Package ‘agricolae’. cran.r-project.org. Disponible en:\url{ https://cran.r-project.org/web/packages/agricolae/agricolae.pdf}. Consultada Enero 2018.\\

\noindent
% Denny,M. (2017). The fallacy of the average: on the ubiquity, utility and continuing novelty of Jensen’s inequality. Journal of Experimental Biology 220, 139-146.\\

\noindent
% Dixon, P. (2003).VEGAN, a package of R functions for community ecology.Journal of Vegetation Science 14: 927-930.\\

\noindent
% Edmonson, R.N.(2014).Agridat.Journal of Agricultural Science , 152, 2.\\

\noindent
%Filippaa,G., Cremonesea,E., Migliavacca, m., Galvagnoa, M., Forkel,M., Wingate,L., Tomelleri,E., Morra,U., Richardsone,A.D. (2016). Phenopix: A R package for image-based vegetation phenology.Agricultural and Forest Meteorology 220, 141-150.\\

\noindent
% Gaitán, Ángela, María Gonzales, Carlos Ñústez, Tatiana Saldaña, and José Miguel Cotes. 2013. “Análisis Funcional de Crecimiento Y Desarrollo de Cuatro Variedades de Papa (Solanum tuberosum Subsp. andigena).” Revista Facultad de Ciencias Básicas 9(2):172-85.\\

\noindent
% Gardner, F., R. Pearce, and R. Mitchell. (1985). Physiology of Crop Plants. Iowa: Iowa State University Press.\\

\noindent
% Grime, J. and R. Hunt. (1975). “Relative Growth-Rate: Its Range and Adaptive Significance in a Local Flora.” Journal of Ecology 63(2):393-422.\\

\noindent
% Hendrik, Poorter. (1989). “Plant Growth Analysis: Towards a Synthesis of the Classical and the Functional Approach.” Physiologia Plantarum 75:237-44.\\

\noindent
% Horton, N., Kleinman, K. (2010). Using R and RStudio for Data Management, Statistical Analysis, and Graphics. CRC Press, A Chapman $\&$ Hall Book. USA. \\

\noindent
% Hunt, R. (1978). “Plant Growth Analysis: The Rationale behind the Use of Fitted Mathematical Function.” Annals of Botany 43:245-49.\\

\noindent
% Hunt, R. (1982). Plant Growth Curves: The Functional Approach to Plant Growth Analysis. New York: Cambridge University Press.\\

\noindent
% Hunt, R. (2002). “A Modern Tool for Classical Plant Growth Analysis.” Annals of Botany 90(4):485–88.\\ 

\noindent
% Hunt, R. (2003). “Growth Analysis, Individual Plants.” 579–88.\\

\noindent
% Kelley, C.T.(1987) Iterative Methods for Optimization, Society for Industrial and Applied Mathematics, Frontiers in Applied Mathematics.\\

\noindent
% Kniss, AR.,  Streibig JC. (2015). Statistical Analysis of Agricultural Experiments with R. (Disponible en:\url{ http://rstats4ag.org/}, Consultado: Enero 2018.\\

\noindent
% Larcher, W. (2003). Physiological Plant Ecology. Ecophysiology and Stress Physiology of the Functional Groups. Fourth. Berlin: Springer.\\

\noindent
%Leisch, F. (2002), “Sweave, Part I: Mixing R and LATEX,” R News, 2, 28–31, URL http://CRAN.R-project.org/doc/Rnews/.\\

\noindent
%Leisch F. (2009). Creating R Packages: A Tutorial.\\

\noindent
%Marsden J., Weinstein A. (1985). Calculus I. Springer-Verlag, New York Inc.\\

\noindent
% Park,B.B., Park,G.E.,Bae,K- (2015).Diagnosis of plant nutrient and growth responses on fertilization with vector analysis and
%morphological index.Forest Science and Technology, Vol. 11, No. 1,1-10.\\

\noindent	
% International Rice Research Institute (IRRI),  Biometrics and Breeding Informatics (BBI) group. Package STAT. Statistical Tool for Agricultural Research. Philippine  Disponible para descarga en:\url{ https://drive.google.com/file/d/0Bw-PTBz1SHmsM0VRTFc0MEZqQVU/view}. Consultado Enero 2018.\\

\noindent	
% Radford, PJ. 1967. “Growth Analysis Formulae - Their Use and Abuse.” Crop Science 7(3):171.\\

\noindent
% Salisbury, F.B. (1996). Units, Symbols, and Terminology for Plant Physiology.A Reference for Presentation of Research Results in the Plant Sciences.New York,Oxford, Oxford University Press.\\

\noindent
% Santos, Marcela. 2010. “Evaluación Del Crecimiento, Desarrollo Y Componentes de Rendimiento de Cuatro Cultivares de Papa Criolla En Dos Localidades Del Departamento de Cundinamarca.” Universidad Nacional de Colombia.\\

\noindent
% Sprouffske, K., Wagner,A. (2016). Growthcurver: an R package for obtaining interpretable metrics from microbial growth curves.Sprouffske and Wagner BMC Bioinformatics (2016) 17:172.\\

\noindent
% The R Foundation for Statistical Computing c/o Institute for Statistics and Mathematics Wirtschaftsuniversität Wien. (2014).  R: Software Development Life Cycle A Description of R’s Development, Testing, Release and Maintenance Processes. Disponible en:\url{ } . Consultado Enero 2018.\\

\noindent 
% Thornley, J., Johnson I.(1990). Plant and Crop Modelling: A Mathematical Approach to Plant and Crop Physiology. The Blackburn Press, New Jersey.\\

\noindent
%Trancón B.,  Carl  W., Bolz F.,  Grelck C. (2012).  The Functional Programming Language R and the Paradigm of Dynamic Scientific Programming. (182-197).\\

\noindent
%Wickham H. (2015). R Packages: Organize, Test, Document, and Share Your Code. 1 Edic. Hadley – Wickham.\\

\noindent
%Xie Y. (2017) . Package ‘formatR’. Disponible en:\url{ https://cran.r-project.org/web/packages/formatR/formatR.pdf}. Consultado Enero 2018.\\

\noindent
%Ablan, M. M\'arquez, R. Rivas, Y. Molina, A y Querales, J. (2011, Noviembre). Una librer\'ia en R para validaci\'on de modelos de simulaci\'on. Revista Ciencia e Ingenier\'ia.” , Volumen:(Edici\'on Especial: “Jornada de Modelado y Simulaci\'on”), pp.117-126.\\

\noindent
%Yeater, K. Duke, S and Riedell, W. (2015, Febrero 25). Multivariate Analysis: Greater Insights into Complex Systems. Agronomy Journal , Volumen(107), pp.799-810\\

\noindent
%Fern\'andez Contreras, M.(2017). Modelo de Forrester para evaluar procesos biol\'ogicos, casos de estudio: crecimiento de hongos en medios l\'iquidos y crecimiento de pl\'antulas de piment\'on Capsicum anuum L. biofertilizadas. (Trabajo Especial de Grado de pregrado).Universidad Nacional Experimental del T\'achira. San Crist\'oba, Estado T\'achira. 










%% Incluir la bibliograf'ia. Mirar el archivo "biblio.bib" para m'as detales
%% y un ejemplo.


\end{document}
